\documentclass[10pt,a4paper]{article}
\usepackage[margin=2.2cm]{geometry}
\usepackage{fontspec}
\setmonofont{DejaVu Sans Mono}[Scale=0.84]
\newfontfamily\unifont{DejaVu Serif}[Scale=0.9]
\usepackage{amsmath,amssymb,mathtools}
\setcounter{MaxMatrixCols}{20}
\usepackage{longtable,booktabs,array}
\usepackage{fancyvrb}
\usepackage[dvipsnames]{xcolor}
\usepackage[colorlinks=true,linkcolor=NavyBlue,urlcolor=NavyBlue]{hyperref}
\usepackage{microtype}
\usepackage{needspace}
\usepackage{tocloft}\setlength{\cftbeforesecskip}{0.45em}
\setlength{\parindent}{0pt}\setlength{\parskip}{0.55em}
\setlength{\emergencystretch}{8em}
\renewcommand{\arraystretch}{1.15}
\allowdisplaybreaks
\begin{document}
\title{The Standard Model on the world brane of 7+1 spacetime, the portals of dirac16complex, and neutrino flavour\\[0.4em]\large The complete Standard Model in the gravitational field of 7+1 spacetime, electroweak symmetry breaking, quark and lepton masses and mixing, dirac16complex at the brane as four Dirac spinors, Kaluza–Klein masses, the portals and the coupled field equations, and neutrino flavour conversion in vacuum, in matter, in the Sun and in supernovae — sections 17–19 of the notebook 7+1.nb}
\author{Repository: github.com/once-ere/7plus1 — prepared for Nash\\ Notebook: 7+1.nb (Wolfram Language 15.0.1), evaluated Tue 27 Oct 2026 16:10:34\\ Document generated by tools/make\_docs.py on 2026-10-27}
\date{}
\maketitle
\begin{abstract}
The 7+1 spacetime of this project has one time coordinate $x_0$, three ordinary space coordinates $x_1,x_2,x_3$ and four auxiliary space coordinates $x_4,\dots,x_7$. This document places the complete Standard Model of particle physics on the 3+1-dimensional WORLD BRANE $x_4=x_5=x_6=x_7=0$, in the gravitational field of the whole spacetime, and couples it to the 16-component field dirac16complex, which lives in the whole spacetime. It states the field content and the Lagrangian, checks anomaly cancellation and electroweak symmetry breaking with current data, decomposes dirac16complex at the brane into four ordinary Dirac spinors, derives its Kaluza–Klein masses, defines three portals and the coupled field equations, and works out neutrino flavour conversion with and without dirac16complex. Sections 17–19 of 7+1.nb contain 24 checks (24 True); the independent Python program tools/ref\_sm.py contains 33 (33 True).
\end{abstract}
\tableofcontents
\Needspace{12\baselineskip}\section*{1. Results at a glance}\addcontentsline{toc}{section}{1. Results at a glance}
\begin{itemize}
\item \textbf{Field content.} Gauge group $SU(3)_C\times SU(2)_L\times U(1)_Y$, three generations including right-handed neutrinos, one Higgs doublet; all gauge anomalies cancel in every generation (section 3).
\item \textbf{Electroweak scale.} $v=(\sqrt2G_F)^{-1/2}=246.2196402$ GeV; tree-level $m_W=gv/2$, $m_Z=\sqrt{g^2+g'^2}\,v/2$, $m_H^2=2\lambda v^2$ (section 5).
\item \textbf{Mixing.} CKM matrix with the Jarlskog invariant $J=3.117\times10^{-5}$; PMNS matrix from NuFit-6.0 (section 6).
\item \textbf{Cosmology on the brane.} The brane Friedmann equation gives the age 13.791 Gyr (section 7).
\item \textbf{dirac16complex at the brane} is four brane Dirac spinors; a mode with auxiliary momentum $\vec q$ has the Kaluza–Klein mass $\sqrt{m^2+|\vec q|^2/b_0^2}$ (sections 8–9).
\item \textbf{Portals} (Yukawa, vector, Higgs) couple dirac16complex to the Standard Model at the brane; the Yukawa portal makes the light neutrino mixing non-unitary (section 10).
\item \textbf{Neutrino flavour.} At DUNE ($L=1300$ km, $E=2.5$ GeV) $P(\nu_\mu\to\nu_e)=0.048837$ in vacuum and 0.072327 in matter (section 11).
\end{itemize}
\Needspace{12\baselineskip}\section*{2. The world brane and the total action}\addcontentsline{toc}{section}{2. The world brane and the total action}
Units $\hbar=c=1$. The auxiliary coordinates form a torus $T^4$: $x_k\simeq x_k+L_c$ ($k=4,\dots,7$; coordinate period $L_c$, proper size $bL_c$). The Standard Model lives on the world brane
\begin{gather*}
W=\{x_4=x_5=x_6=x_7=0\},\qquad x=(x_0,x_1,x_2,x_3),\quad y=(x_4,\dots,x_7),\\ h_{\mu\nu}(x)=g_{\mu\nu}(x,y=0)\ \ (\mu,\nu=0,\dots,3),
\end{gather*}
with the induced metric $h_{\mu\nu}$; gravity and dirac16complex live in the whole 7+1 spacetime (the bulk).

\textbf{The signature: $\{7,1,0\}$ and $\mathrm{diag}(+1,-1,\dots,-1)$.} The 7+1 spacetime has one time-like dimension ($x_0$, the direction of the octonion unit $e_0=1$) and seven space-like dimensions ($x_1,x_2,x_3$ ordinary space, $x_4,\dots,x_7$ auxiliary space). Two different objects describe this, and neither may be replaced by the other. (i) The tensor package xAct, used in section 12 of 7+1.nb exactly as in the author’s notebook Einstein-Lovelock-7+1-Nash.nb, declares the metric with \texttt{DefMetric[\{7, 1, 0\}, g71[-a, -b], CD, …]}: the list $\{7,1,0\}$ declares seven space-like dimensions, one time-like dimension and no degenerate dimension — the numbers of entries $+1$ and $-1$ of $\eta_{71}=\mathrm{diag}(-1,+1,\dots,+1)$ (\texttt{η71}), the same spacetime written in the sign convention of 4+4.nb (space-like $=+1$, time-like $=-1$). (ii) The matrix $\eta=-\eta_{71}=\mathrm{diag}(+1,-1,\dots,-1)$ (\texttt{η17}) is simply the squares of the octonion basis: $e_0^2=+1$ for the unit $e_0=1$ (time) and $e_i^2=-1$ for the seven imaginary units $e_1,\dots,e_7$ (space); more generally it is the real part of the product of two basis units, and it gives the quadratic form of the octonions:
\begin{gather*}
\eta_{AB}=\mathrm{Re}(e_Ae_B),\qquad e_0^2=+1,\qquad e_i^2=-1\ \ (i=1,\dots,7),\\ \mathrm{Re}(x\,x)=x_0^2-x_1^2-\cdots-x_7^2=\eta_{AB}\,x^Ax^B\qquad (x=x^Ae_A\in\mathbb O).
\end{gather*}
$\eta$ gives the signs of the Clifford relation $\Gamma_A\Gamma_B+\Gamma_B\Gamma_A=2\,\eta_{AB}\,1_{16}$ of the real 16×16 Dirac matrices: time squares to $+1_{16}$ and each of the seven space directions to $-1_{16}$, exactly like the octonion units (in the notation $Cl(p,q)$, with $p$ generators squaring to $+1$ and $q$ to $-1$, the algebra $Cl(1,7)\cong\mathbb R(16)$). It also gives the signs of the metric components: $g_{00}>0$ along the time-like direction, negative squares along the space-like directions. Every component computed by 7+1.nb follows from the component matrix of the metric; the declared list fixes for xAct the sign of the determinant of the metric, \texttt{SignDetOfMetric[g71]} $=-1$, which is the same in both sign conventions, since $\det \eta=\det \eta_{71}=-1$. 7+1.nb checks \texttt{SignatureOfMetric[g71]} $=\{7,1,0\}$, \texttt{SignDetOfMetric[g71]} $=-1$, $\mathrm{Re}(e_Ae_B)=\eta_{AB}$ and $\mathrm{Re}(x\,x)=\eta_{AB}x^Ax^B$.

The brane has one time-like dimension ($x_0$) and three space-like dimensions ($x_1,x_2,x_3$) — in the counting of xAct the list $\{3,1,0\}$ — and its flat frame metric $\mathrm{diag}(+1,-1,-1,-1)$ is $\eta$ restricted to $x_0,\dots,x_3$: the squares of the basis $e_0=1,e_1,e_2,e_3$ of the quaternion subalgebra $\mathbb H$ of the octonions (checked in section 17 of 7+1.nb). The total action is
\begin{equation*}
S=S_{\rm EL}[g]+S_{\psi}[\psi,g]+\int_Wd^4x\,\sqrt{|h|}\,\big(\mathcal L_{\rm SM}+\mathcal L_{\rm portal}\big)+S_{\rm stab},\qquad \int_Wd^4x\,\sqrt{|h|}\,(\dots)=\int d^8X\,\sqrt{|g|}\,\frac{\delta^4(y)}{b^4}\,(\dots),
\end{equation*}
where the last identity holds for the cosmological metric $ds^2=dx_0^2-a^2\sum_{i=1}^3dx_i^2-b^2\sum_{k=4}^7dx_k^2$ ($\sqrt{|g|}=a^3b^4$, $\sqrt{|h|}=a^3$). The pieces:
\begin{align*}
S_{\rm EL}&=-\frac1{2\kappa}\int d^8X\sqrt{|g|}\sum_{k=0}^{3}\alpha_kL_k,\qquad L_k=\frac1{2^k}\delta^{\alpha_1\beta_1\cdots\alpha_k\beta_k}_{\gamma_1\delta_1\cdots\gamma_k\delta_k}R^{\gamma_1\delta_1}{}_{\alpha_1\beta_1}\cdots R^{\gamma_k\delta_k}{}_{\alpha_k\beta_k}, \\
S_\psi&=\int d^8X\sqrt{|g|}\Big[\bar\psi\big(i\Gamma^\mu(X)D_\mu-m\big)\psi-\frac12\sum_{r=0}^{8}g_r\sum_{|K|=r}(\bar\psi\Theta_K\psi)(\bar\psi\Theta^K\psi)\Big],
\end{align*}
($\kappa=8\pi G_8$; $\alpha_0=2\bar\Lambda$, $\alpha_1=1$; $\Gamma^\mu(X)$ the generalized frame matrices, $D_\mu=\partial_\mu+\Omega_\mu$ with the canonical spin connection $\Omega_\mu$; $\Theta^K$ the 256 Clifford products with $g_r$ the nine self-couplings of dirac16complex), and $S_{\rm stab}$ is a stress that keeps $b$ constant in the modern epoch.

\Needspace{12\baselineskip}\section*{3. Fields, quantum numbers and anomaly cancellation (section 17.1)}\addcontentsline{toc}{section}{3. Fields, quantum numbers and anomaly cancellation (section 17.1)}
Gauge fields $G^a_\mu$ ($a=1,\dots,8$), $W^I_\mu$ ($I=1,2,3$), $B_\mu$ with couplings $g_s$, $g$, $g'$. Matter in each of three generations (Weyl fields; electric charge $Q=T_3+Y$):

{\small
\begin{longtable}{>{\raggedright\arraybackslash}p{\dimexpr 0.2000\linewidth-2\tabcolsep\relax}>{\raggedright\arraybackslash}p{\dimexpr 0.1200\linewidth-2\tabcolsep\relax}>{\raggedright\arraybackslash}p{\dimexpr 0.1200\linewidth-2\tabcolsep\relax}>{\raggedright\arraybackslash}p{\dimexpr 0.1000\linewidth-2\tabcolsep\relax}>{\raggedright\arraybackslash}p{\dimexpr 0.1800\linewidth-2\tabcolsep\relax}>{\raggedright\arraybackslash}p{\dimexpr 0.2800\linewidth-2\tabcolsep\relax}}
\toprule
\textbf{field} & \textbf{$SU(3)_C$} & \textbf{$SU(2)_L$} & \textbf{$Y$} & \textbf{chirality $\chi_f$} & \textbf{charges} \\
\midrule\endfirsthead
\toprule
\textbf{field} & \textbf{$SU(3)_C$} & \textbf{$SU(2)_L$} & \textbf{$Y$} & \textbf{chirality $\chi_f$} & \textbf{charges} \\
\midrule\endhead
\endfoot
\bottomrule\endlastfoot
$Q_L=(u_L,d_L)$ & 3 & 2 & $1/6$ & $+1$ (left) & $2/3$, $-1/3$ \\
$u_R$ & 3 & 1 & $2/3$ & $-1$ (right) & $2/3$ \\
$d_R$ & 3 & 1 & $-1/3$ & $-1$ & $-1/3$ \\
$L_L=(\nu_L,e_L)$ & 1 & 2 & $-1/2$ & $+1$ & $0$, $-1$ \\
$e_R$ & 1 & 1 & $-1$ & $-1$ & $-1$ \\
$\nu_R$ & 1 & 1 & $0$ & $-1$ & $0$ \\
$\Phi$ (Higgs) & 1 & 2 & $1/2$ & — & $\tilde\Phi=i\sigma_2\Phi^*$ \\*
\end{longtable}}
With $N_f$ the number of colours and $d_f$ the $SU(2)$ dimension of field $f$, the anomaly coefficients of one generation are
\begin{align*}
A_{SU(3)^2U(1)}&=\sum_{f:\,N_f=3}\chi_fd_fY_f,\qquad A_{SU(2)^2U(1)}=\sum_{f:\,d_f=2}\chi_fN_fY_f,\qquad A_{SU(3)^3}=\sum_{f:\,N_f=3}\chi_fd_f, \\
A_{U(1)^3}&=\sum_f\chi_fN_fd_fY_f^3,\qquad A_{\rm grav^2U(1)}=\sum_f\chi_fN_fd_fY_f,
\end{align*}
and all five vanish, with or without $\nu_R$; the number of $SU(2)$ doublets per generation ($3+1=4$) is even (no Witten anomaly). The relativistic degrees of freedom above the electroweak scale (without $\nu_R$) are
\begin{equation*}
g_*=\underbrace{16}_{G}+\underbrace{6}_{W}+\underbrace{2}_{B}+\underbrace{4}_{\Phi}+\frac78\cdot3\cdot\big(12+6+6+4+2\big)=28+\frac78\cdot90=106.75 .
\end{equation*}
\Needspace{12\baselineskip}\section*{4. The Lagrangian on the brane, in the gravitational field (section 17.2)}\addcontentsline{toc}{section}{4. The Lagrangian on the brane, in the gravitational field (section 17.2)}
With the frame (tetrad) of $h_{\mu\nu}$, its spin connection and complex 4×4 Dirac matrices $\gamma^\mu(x)$ of the brane ($\gamma^\mu\gamma^\nu+\gamma^\nu\gamma^\mu=2h^{\mu\nu}$), and the gauge-covariant derivative
\begin{equation*}
D_\mu=\nabla_\mu-ig_s\,G^a_\mu T^a-ig\,W^I_\mu\frac{\tau^I}2-ig'\,Y\,B_\mu
\end{equation*}
($T^a$ the generators of $SU(3)$ in the representation of the field, $\tau^I$ the Pauli matrices; $\nabla_\mu$ includes the spin connection of $h$ on spinors), the Standard Model Lagrangian is
\begin{align*}
\mathcal L_{\rm SM}&=-\tfrac14G^a_{\mu\nu}G^{a\mu\nu}-\tfrac14W^I_{\mu\nu}W^{I\mu\nu}-\tfrac14B_{\mu\nu}B^{\mu\nu}+(D_\mu\Phi)^\dagger(D^\mu\Phi)-\mu^2\Phi^\dagger\Phi-\lambda(\Phi^\dagger\Phi)^2 \\
&\quad+\sum_f\bar f\,i\gamma^\mu D_\mu f-\big[\bar Q_LY_u\tilde\Phi u_R+\bar Q_LY_d\Phi d_R+\bar L_LY_e\Phi e_R+\bar L_LY_\nu\tilde\Phi\nu_R+\mathrm{h.c.}\big], \\
G^a_{\mu\nu}&=\partial_\mu G^a_\nu-\partial_\nu G^a_\mu+g_sf^{abc}G^b_\mu G^c_\nu,\quad W^I_{\mu\nu}=\partial_\mu W^I_\nu-\partial_\nu W^I_\mu+g\,\epsilon^{IJK}W^J_\mu W^K_\nu,\quad B_{\mu\nu}=\partial_\mu B_\nu-\partial_\nu B_\mu,
\end{align*}
$f$ runs over all Weyl fields of the three generations, and $Y_u,Y_d,Y_e,Y_\nu$ are complex 3×3 Yukawa matrices. Its field equations are the ordinary ones with $\partial\to\nabla$ (the Levi-Civita connection of $h$), e.g. $\nabla_\mu G^{a\mu\nu}+g_sf^{abc}G^b_\mu G^{c\mu\nu}=-g_s\sum_q\bar q\gamma^\nu T^aq$.

\Needspace{12\baselineskip}\section*{5. Electroweak symmetry breaking (section 17.3)}\addcontentsline{toc}{section}{5. Electroweak symmetry breaking (section 17.3)}
With $\langle\Phi\rangle=(0,v/\sqrt2)$ the term $(D_\mu\langle\Phi\rangle)^\dagger(D^\mu\langle\Phi\rangle)$ is a mass matrix for $(W^1,W^2,W^3,B)$, its Hessian $M^2_{ij}=\partial^2\big|g\,W^I\tfrac{\tau^I}2\langle\Phi\rangle+\tfrac{g'}2B\langle\Phi\rangle\big|^2/\partial w_i\partial w_j$ having the eigenvalues
\begin{equation*}
m_W^2=\frac{g^2v^2}4\ \ (\text{twice}),\qquad m_Z^2=\frac{(g^2+g'^2)v^2}4,\qquad m_\gamma^2=0 .
\end{equation*}
Inputs (Particle Data Group 2025): $G_F=1.1663788\times10^{-5}$ GeV$^{-2}$, $m_W=80.3692$ GeV, $m_Z=91.1880$ GeV, $m_H=125.20$ GeV. At tree level (on-shell weak angle)
\begin{equation*}
v=(\sqrt2G_F)^{-1/2},\quad g=\frac{2m_W}{v},\quad \cos\theta_W=\frac{m_W}{m_Z},\quad g'=g\tan\theta_W,\quad e=g\sin\theta_W,\quad \lambda=\frac{m_H^2}{2v^2},
\end{equation*}
{\footnotesize
\begin{longtable}{>{\raggedright\arraybackslash}p{\dimexpr 0.1500\linewidth-2\tabcolsep\relax}>{\raggedright\arraybackslash}p{\dimexpr 0.1400\linewidth-2\tabcolsep\relax}>{\raggedright\arraybackslash}p{\dimexpr 0.1400\linewidth-2\tabcolsep\relax}>{\raggedright\arraybackslash}p{\dimexpr 0.1500\linewidth-2\tabcolsep\relax}>{\raggedright\arraybackslash}p{\dimexpr 0.1400\linewidth-2\tabcolsep\relax}>{\raggedright\arraybackslash}p{\dimexpr 0.1400\linewidth-2\tabcolsep\relax}>{\raggedright\arraybackslash}p{\dimexpr 0.1400\linewidth-2\tabcolsep\relax}}
\toprule
\textbf{$v$ (GeV)} & \textbf{$g$} & \textbf{$g'$} & \textbf{$\sin^2\theta_W$} & \textbf{$e$} & \textbf{$1/\alpha=4\pi/e^2$} & \textbf{$\lambda$} \\
\midrule\endfirsthead
\toprule
\textbf{$v$ (GeV)} & \textbf{$g$} & \textbf{$g'$} & \textbf{$\sin^2\theta_W$} & \textbf{$e$} & \textbf{$1/\alpha=4\pi/e^2$} & \textbf{$\lambda$} \\
\midrule\endhead
\endfoot
\bottomrule\endlastfoot
246.2196402 & 0.6528252573 & 0.3499462340 & 0.2232094928 & 0.3084276316 & 132.1001486 & 0.1292805757 \\*
\end{longtable}}
($1/\alpha$ here is the tree-level value at the weak scale, not the low-energy fine-structure constant.)

\Needspace{12\baselineskip}\section*{6. Quark and lepton masses; CKM and PMNS mixing (section 17.4)}\addcontentsline{toc}{section}{6. Quark and lepton masses; CKM and PMNS mixing (section 17.4)}
Fermion masses are $m=yv/\sqrt2$ for the singular values $y$ of the Yukawa matrices; the mixing matrices are $V_{\rm CKM}=U_{uL}^\dagger U_{dL}$ and $U_{\rm PMNS}=U_{eL}^\dagger U_{\nu L}$, where $Y=U_L\,\mathrm{diag}(y)\,U_R^\dagger$. Both are written in the standard parametrization ($c_{ij}=\sqrt{1-s_{ij}^2}$):
\begin{equation*}
U(s_{12},s_{13},s_{23},\delta)=\begin{pmatrix}1&0&0\\0&c_{23}&s_{23}\\0&-s_{23}&c_{23}\end{pmatrix}\begin{pmatrix}c_{13}&0&s_{13}e^{-i\delta}\\0&1&0\\-s_{13}e^{i\delta}&0&c_{13}\end{pmatrix}\begin{pmatrix}c_{12}&s_{12}&0\\-s_{12}&c_{12}&0\\0&0&1\end{pmatrix}.
\end{equation*}
\textbf{Quarks} (PDG 2024 CKM fit: $s_{12}=0.22501$, $s_{13}=0.003732$, $s_{23}=0.04183$, $\delta=1.147$; PDG 2025 masses: $m_u=2.16$ MeV, $m_d=4.70$ MeV, $m_s=93.5$ MeV, $m_c=1.2730$ GeV, $m_b=4.183$ GeV, $m_t=172.56$ GeV). With $Y_u=\sqrt2\,\mathrm{diag}(m_u,m_c,m_t)/v$ and $Y_d=\sqrt2\,V\,\mathrm{diag}(m_d,m_s,m_b)/v$ the singular values reproduce the masses and $|U_{uL}^\dagger U_{dL}|=|V|$. The notebook obtains
\begin{equation*}
|V_{\rm CKM}|=\begin{pmatrix}0.974350&0.225008&0.00373200\\0.224876&0.973489&0.0418297\\0.00858280&0.0411095&0.999118\end{pmatrix},
\end{equation*}
\begin{gather*}
J=\mathrm{Im}\big(V_{us}V_{cb}V_{ub}^*V_{cs}^*\big)=3.11699\times10^{-5},\qquad \bar\rho+i\bar\eta=-\frac{V_{ud}V_{ub}^*}{V_{cd}V_{cb}^*}=0.159184+0.352275\,i,\\ A=\frac{s_{23}}{s_{12}^2}=0.826198,\qquad y_t=\frac{\sqrt2m_t}{v}=0.99113414,
\end{gather*}
in agreement with the PDG 2024 values $J=(3.12^{+0.13}_{-0.12})\times10^{-5}$, $A=0.826$, $\bar\rho=0.1591$, $\bar\eta=0.3523$.

\textbf{Leptons} (NuFit-6.0, normal ordering, without the atmospheric data of Super-Kamiokande: $\sin^2\theta_{12}=0.307$, $\sin^2\theta_{23}=0.561$, $\sin^2\theta_{13}=0.02195$, $\delta=177^\circ$, $\Delta m^2_{21}=7.49\times10^{-5}$ eV², $\Delta m^2_{31}=2.534\times10^{-3}$ eV², $\Delta m^2_{ij}=m_i^2-m_j^2$). With $\nu_R$ the neutrinos are Dirac particles; with the lightest mass $m_1=0$: $m_2=\sqrt{\Delta m^2_{21}}$, $m_3=\sqrt{\Delta m^2_{31}}$, and $Y_\nu=\sqrt2\,U\,\mathrm{diag}(m_1,m_2,m_3)/v$ (electron Yukawa basis diagonal):
\begin{gather*}
|U_{\rm PMNS}|=\begin{pmatrix}0.82328&0.54796&0.14816\\0.27491&0.61298&0.74073\\0.49663&0.56921&0.65526\end{pmatrix},\qquad J_{\rm PMNS}=0.0017359,\\ (m_1,m_2,m_3)=(0, 0.00865448, 0.0503389)\ \text{eV}.
\end{gather*}
\Needspace{12\baselineskip}\section*{7. The modern epoch on the brane (section 17.5)}\addcontentsline{toc}{section}{7. The modern epoch on the brane (section 17.5)}
With $b=b_0$ the 00 field equation is the Friedmann equation of the brane,
\begin{equation*}
3H_a^2=8\pi G_4\,\rho_4+\Lambda_4,\qquad G_4=\frac{G_8}{V_\perp},\qquad V_\perp=b_0^4L_c^4,\qquad \varepsilon_8=\frac{\rho_4}{V_\perp}\ \ \text{(matter on the brane)},
\end{equation*}
($H_a=a'/a$, $\rho_4$ the four-dimensional energy density). With the Planck 2018 parameters $H_0=67.4$ km/s/Mpc, $\Omega_m=0.315$, $\Omega_r=9.14\times10^{-5}$, $\Omega_\Lambda=1-\Omega_m-\Omega_r$ the age of the universe is
\begin{equation*}
t_0=\frac1{H_0}\int_0^1\frac{da}{a\sqrt{\Omega_r a^{-4}+\Omega_m a^{-3}+\Omega_\Lambda}}=13.7907\ \text{Gyr}\qquad(\text{Planck 2018: }13.787\pm0.020).
\end{equation*}
With the reduced eight-dimensional Planck mass $M_8$ ($8\pi G_8=M_8^{-6}$) the reduced four-dimensional one is $M_{\rm Pl}^2=M_8^6V_\perp$ ($M_{\rm Pl}=2.435\times10^{18}$ GeV). For $M_8=10$ TeV:
\begin{equation*}
V_\perp=\frac{M_{\rm Pl}^2}{M_8^6}=5.92923\times10^{12}\ \text{GeV}^{-4},\qquad V_\perp^{1/4}=3.07919\times10^{-13}\ \text{m},\qquad \frac{2\pi}{V_\perp^{1/4}}=4.02652\ \text{MeV},
\end{equation*}
the last being the lowest Kaluza–Klein mass of a field on the auxiliary torus (proper period $V_\perp^{1/4}=b_0L_c$).

\Needspace{12\baselineskip}\section*{8. dirac16complex at the brane: four brane Dirac spinors (section 18)}\addcontentsline{toc}{section}{8. dirac16complex at the brane: four brane Dirac spinors (section 18)}
At the brane the eight Dirac matrices split into four along $W$ and four transverse to it. With the complex 4×4 Dirac matrices of the Dirac representation and four anti-Hermitian 4×4 matrices $\rho_k$,
\begin{align*}
\gamma_0&=\mathrm{diag}(1,1,-1,-1),\quad \gamma_i=\begin{pmatrix}0&\sigma_i\\-\sigma_i&0\end{pmatrix},\quad \gamma_5=i\gamma_0\gamma_1\gamma_2\gamma_3,\qquad \gamma_\mu\gamma_\nu+\gamma_\nu\gamma_\mu=2\eta_{\mu\nu}1_4, \\
\rho_4&=i\gamma_0,\quad \rho_5=-\gamma_1,\quad \rho_6=-\gamma_2,\quad \rho_7=-\gamma_3,\qquad \rho_k\rho_l+\rho_l\rho_k=-2\delta_{kl}1_4,
\end{align*}
($\sigma_i$ the Pauli matrices, $\eta=\mathrm{diag}(1,-1,-1,-1)$, the squares of the quaternion basis $1,e_1,e_2,e_3$) the sixteen-dimensional matrices
\begin{equation*}
R_\mu=\gamma_\mu\otimes1_4\ \ (\mu=0,\dots,3),\qquad R_k=\gamma_5\otimes\rho_k\ \ (k=4,\dots,7)
\end{equation*}
satisfy the same Clifford relations as the real Dirac matrices $\Gamma_A$ of 7+1 spacetime, $R_AR_B+R_BR_A=2\eta_{AB}1_{16}$ with $\eta=\mathrm{diag}(+1,-1,\dots,-1)$, the squares of the octonion basis. Both representations are irreducible, so there is a unitary $U$ with $\Gamma_AU=UR_A$; it is obtained as the Schur average over the 256 Clifford products,
\begin{equation*}
U\propto\sum_{K}\Gamma^{(K)}\,s\,\big(R^{(K)}\big)^{-1},\qquad \Gamma^{(K)}=\Gamma_{A_1}\cdots\Gamma_{A_r},\quad R^{(K)}=R_{A_1}\cdots R_{A_r},
\end{equation*}
with a fixed matrix $s$ (the first matrix unit for which the sum is non-zero), normalized so that $U^\dagger U=1_{16}$. Writing $\psi=U\zeta$ with $\zeta\in\mathbb C^4\otimes\mathbb C^4$, $\zeta=\sum_{a=1}^4\psi_a\otimes e_a$, dirac16complex at the brane is FOUR brane Dirac spinors $\psi_a$ ($a$ the transverse spinor index), and with $\bar\psi_a=\psi_a^\dagger\gamma_0$, $\gamma^\mu=\eta^{\mu\mu}\gamma_\mu$:
\begin{equation*}
\bar\psi\psi'=\sum_{a=1}^4\bar\psi_a\psi'_a,\qquad \bar\psi\,\Gamma^\mu\psi'=\sum_{a=1}^4\bar\psi_a\gamma^\mu\psi'_a\ \ (\mu=0,\dots,3).
\end{equation*}
The eight-dimensional chirality $\chi=i\Gamma_0\cdots\Gamma_7$ is
\begin{equation*}
\chi=U\,(\gamma_5\otimes\rho_{4567})\,U^\dagger,\qquad \rho_{4567}=\rho_4\rho_5\rho_6\rho_7\ \ \text{(Hermitian, eigenvalues }+1,+1,-1,-1),
\end{equation*}
so the left-handed half $\psi_L$ ($\chi=-1$) contains brane-right-handed parts with $\rho_{4567}=-1$ and brane-left-handed parts with $\rho_{4567}=+1$.

\Needspace{12\baselineskip}\section*{9. Kaluza–Klein masses on the auxiliary torus (section 18.1)}\addcontentsline{toc}{section}{9. Kaluza–Klein masses on the auxiliary torus (section 18.1)}
With $b=b_0$ and $\psi\propto e^{i\vec q\cdot\vec y}$, $\vec q=2\pi\vec n/L_c$ ($\vec n\in\mathbb Z^4$), the eight-dimensional Dirac equation at the brane becomes, in the variables $\zeta$,
\begin{equation*}
\big(\gamma\cdot p\otimes1_4+X-m\big)\zeta=0,\qquad \gamma\cdot p=p_0\gamma_0-\sum_{i=1}^3p_i\gamma_i,\qquad X=\gamma_5\otimes\frac{\vec\rho\cdot\vec q}{b_0},\qquad X^2=-\frac{|\vec q|^2}{b_0^2}1_{16},
\end{equation*}
and $X$ anticommutes with every $\gamma_\mu\otimes1_4$. Hence each Fourier mode is four brane Dirac fermions of mass
\begin{equation*}
m_{\vec n}=\sqrt{m^2+\frac{|\vec q|^2}{b_0^2}}\qquad(\text{the zero mode }\vec n=0\text{ has mass }m),
\end{equation*}
checked by the vanishing of $\det(\gamma\cdot p\otimes1_4+X-m)$ on $p_0^2-|\vec p|^2=m_{\vec n}^2$ (and not off shell), in 40-digit arithmetic, for $m=7/10$, $b_0=13/10$, $\vec q=2\pi(2,-1,\frac12,3)$, $\vec p=(0.3,-0.2,0.4)$.

\Needspace{12\baselineskip}\section*{10. The portals and the coupled field equations (section 18.2)}\addcontentsline{toc}{section}{10. The portals and the coupled field equations (section 18.2)}
The portals of dirac16complex are local couplings at the brane (an invented interaction; separate couplings for the two eight-dimensional chiral halves), with $\psi_a(x)$ the brane spinors of $\psi(x,y=0)$:

\begin{itemize}
\item \textbf{Neutrino (Yukawa) portal:} $-\sum_{i,s}\big[y^L_{is}\,\bar L_i\tilde\Phi P_R\psi^{(L)}_s+y^R_{is}\,\bar L_i\tilde\Phi P_R\psi^{(R)}_s\big]+\mathrm{h.c.}$, where $\psi^{(L)}_s$ ($\psi^{(R)}_s$), $s=1,2$, are the brane-right-handed parts of the left (right) half of dirac16complex and $P_R=\frac12(1+\gamma_5)$. After symmetry breaking it is a Dirac mass $\mu_{ia}=y_{ia}v/(\sqrt2\sqrt{V_\perp})$ between $\nu_{L,i}$ and every Kaluza–Klein mode ($1/\sqrt{V_\perp}$ normalizes the modes on the torus).
\item \textbf{Vector (current) portal:} $-G^f_P\,(\bar\psi\Gamma^\mu\psi)|_W\,(\bar f\gamma_\mu f)$ for Standard Model fermions $f$. In a background of dirac16complex with brane density $n_\psi$ it adds the potential $V^f_P=G^f_Pn_\psi$ to the forward scattering of $f$ (for neutrinos: a non-standard matter effect).
\item \textbf{Higgs (scalar) portal:} $-\lambda_P\,\Phi^\dagger\Phi\,(\bar\psi\psi)|_W$, which shifts the dirac16complex mass by $\lambda_Pv^2/(2V_\perp)$ after symmetry breaking.
\end{itemize}
The coupled field equations of all fields:
\begin{align*}
\sum_k\alpha_kE^{(k)\mu}{}_\nu&=\kappa\Big(T_\psi+T_{\rm stab}+\frac{\delta^4(y)}{b^4}T_W\Big)^\mu{}_\nu, \\
\big(i\Gamma^\mu D_\mu-m\big)\psi-\sum_rg_r\sum_{|K|=r}(\bar\psi\Theta_K\psi)\Theta^K\psi&=\frac{\delta^4(y)}{b^4}\,J_{\rm portal},
\end{align*}
with $E^{(k)}$ the Lovelock tensors, $T_W$ the energy–momentum tensor of the Standard Model fields and portals on the brane, $J_{\rm portal}=\partial\mathcal L_{\rm portal}/\partial\bar\psi$ (for the Yukawa portal $U\big(\sum_{i}y^*_{is}\tilde\Phi^\dagger L_i\otimes e_s+\dots\big)$), and the Standard Model field equations of section 4 with the portal terms as additional sources. In the neutrino sector the Yukawa portal gives, in the basis $(\nu_L,\psi_L)\times(\nu_R,\psi_R)$, the Dirac mass matrix of the zero mode
\begin{equation*}
\mathcal M=\begin{pmatrix}m_D&\mu\\0&m\,1_4\end{pmatrix},\qquad m_D=U_{\rm PMNS}\,\mathrm{diag}(m_1,m_2,m_3),
\end{equation*}
and for $m\gg\mu$ the light mixing matrix $N=(1-\eta)U_{\rm PMNS}$ is no longer unitary: $1-NN^\dagger=2\eta=\mu\mu^\dagger/m^2$ to second order, while the heavy states are mostly dirac16complex. Example: $m=1$ GeV and $\mu$ (in eV)
\begin{equation*}
\mu=\begin{pmatrix}3&1&0&0\\1&2&0&0\\-2&1&0&0\end{pmatrix}\times10^{4}:
\end{equation*}
\begin{gather*}
\text{light masses}=(0, 0.0086544743, 0.050338833)\ \text{eV},\\ \text{heavy masses}=(1.00000000000, 1.00000000000, 1.00000000025, 1.00000000075)\ \text{GeV},\\ (1-NN^\dagger)_{ee}=1.00000\times10^{-9} .
\end{gather*}
\Needspace{12\baselineskip}\section*{11. Neutrino flavour and flavour conversion (section 19)}\addcontentsline{toc}{section}{11. Neutrino flavour and flavour conversion (section 19)}
A neutrino of flavour $\alpha$ produced with energy $E$ is the superposition $\sum_iU^*_{\alpha i}|\nu_i\rangle$. In vacuum (with $\hbar c=1.973269804\times10^{-16}$ GeV m) the amplitude after a distance $L$ is $A(\alpha\to\beta)=\sum_iU_{\beta i}U^*_{\alpha i}e^{-im_i^2L/(2E)}$, so
\begin{equation*}
P(\alpha\to\beta)=\delta_{\alpha\beta}-4\sum_{i>j}\mathrm{Re}\,X_{ij}\sin^2\Delta_{ij}+2\sum_{i>j}\mathrm{Im}\,X_{ij}\sin2\Delta_{ij},\qquad X_{ij}=U^*_{\alpha i}U_{\beta i}U_{\alpha j}U^*_{\beta j},\quad \Delta_{ij}=\frac{\Delta m^2_{ij}L}{4E}.
\end{equation*}
In matter of electron density $n_e$ the charged-current forward scattering adds $V_{CC}=\sqrt2G_Fn_e$ to the $\nu_e$ entry of the flavour Hamiltonian ($V_{CC}\to-V_{CC}$ and $U\to U^*$ for antineutrinos; a potential common to all flavours drops out):
\begin{equation*}
H=\frac1{2E}\,U\,\mathrm{diag}(0,\Delta m^2_{21},\Delta m^2_{31})\,U^\dagger+\mathrm{diag}(V_{CC},0,0),\qquad V_{CC}\big(\rho Y_e=1\ \text{g/cm}^3\big)=7.63\times10^{-14}\ \text{eV}.
\end{equation*}
In two flavours the mixing angle in matter obeys
\begin{equation*}
\tan2\theta_m=\frac{\sin2\theta}{\cos2\theta-2EV/\Delta m^2},
\end{equation*}
maximal at the MSW resonance $2EV=\Delta m^2\cos2\theta$. \textbf{Long baseline} (DUNE: $L=1300$ km, $E=2.5$ GeV, $\rho=2.848$ g/cm³, $Y_e=0.5$):

{\small
\begin{longtable}{>{\raggedright\arraybackslash}p{\dimexpr 0.3400\linewidth-2\tabcolsep\relax}>{\raggedright\arraybackslash}p{\dimexpr 0.2200\linewidth-2\tabcolsep\relax}>{\raggedright\arraybackslash}p{\dimexpr 0.2200\linewidth-2\tabcolsep\relax}>{\raggedright\arraybackslash}p{\dimexpr 0.2200\linewidth-2\tabcolsep\relax}}
\toprule
\textbf{} & \textbf{$P(\nu_\mu\to\nu_e)$} & \textbf{$P(\bar\nu_\mu\to\bar\nu_e)$} & \textbf{$P(\nu_\mu\to\nu_\mu)$} \\
\midrule\endfirsthead
\toprule
\textbf{} & \textbf{$P(\nu_\mu\to\nu_e)$} & \textbf{$P(\bar\nu_\mu\to\bar\nu_e)$} & \textbf{$P(\nu_\mu\to\nu_\mu)$} \\
\midrule\endhead
\endfoot
\bottomrule\endlastfoot
vacuum & 0.048836853 & 0.050198510 & 0.012813139 \\
matter & 0.072326938 & 0.029191799 & — \\
matter + vector portal $V_{P,e}=0.2V_{CC}$ & 0.077049983 & — & — \\*
\end{longtable}}
\textbf{The Sun.} With the electron density $n_e(r)=245\,N_A\,e^{-10.54\,r/R_\odot}$ cm$^{-3}$ ($R_\odot=6.957\times10^8$ m, $N_A$ Avogadro’s number) the two-flavour evolution (effective potential $c_{13}^2V_{CC}$) is integrated from the centre to the surface; the result agrees with the adiabatic MSW formula
\begin{equation*}
P_{ee}=c_{13}^4\,\frac{1+\cos2\theta_{12}\cos2\theta_m(0)}2+s_{13}^4,\qquad \text{vacuum average: }c_{13}^4\Big(1-\frac{\sin^22\theta_{12}}2\Big)+s_{13}^4,
\end{equation*}
{\small
\begin{longtable}{>{\raggedright\arraybackslash}p{\dimexpr 0.2000\linewidth-2\tabcolsep\relax}>{\raggedright\arraybackslash}p{\dimexpr 0.4000\linewidth-2\tabcolsep\relax}>{\raggedright\arraybackslash}p{\dimexpr 0.4000\linewidth-2\tabcolsep\relax}}
\toprule
\textbf{$E$ (MeV)} & \textbf{$P_{ee}$ (numerical)} & \textbf{$P_{ee}$ (adiabatic formula)} \\
\midrule\endfirsthead
\toprule
\textbf{$E$ (MeV)} & \textbf{$P_{ee}$ (numerical)} & \textbf{$P_{ee}$ (adiabatic formula)} \\
\midrule\endhead
\endfoot
\bottomrule\endlastfoot
0.300 & 0.52516444 & 0.52516451 \\
1.00 & 0.45841194 & 0.45841225 \\
5.00 & 0.31033778 & 0.31033351 \\
10.0 & 0.29792183 & 0.29791749 \\*
\end{longtable}}
(at 0.3 MeV the residual matter effect at the centre, $2EV/\Delta m^2_{21}\approx0.15$, still lowers $P_{ee}$ below the vacuum average by about 0.025).

\textbf{Supernovae.} The MSW resonance densities $\rho_{\rm res}=\Delta m^2\cos2\theta/(2EV_{CC}(1\,\text{g/cm}^3))$ for $E=20$ MeV, $Y_e=0.5$ are 1587.14 g/cm³ (H: $\Delta m^2_{31}$, $\theta_{13}$) and 18.9397 g/cm³ (L: $\Delta m^2_{21}$, $\theta_{12}$).

\textbf{Cosmic neutrino background.} $T_\nu=(4/11)^{1/3}T_\gamma$ with $T_\gamma=2.7255$ K gives $T_\nu=1.94537$ K and the number density per flavour (neutrinos and antineutrinos) $n_\nu=\frac{3\zeta(3)}{2\pi^2}T_\nu^3=112.016$ cm$^{-3}$.

\textbf{dirac16complex and neutrino flavour.} With the Yukawa portal the oscillation probabilities of the light states, computed with $N$ in place of $U$, have $\sum_\beta P(\alpha\to\beta)<1$ (a zero-distance effect, compensated by the heavy states); for $\mu\to0$ they become the three-flavour vacuum probabilities above. The vector portal acts as a non-standard matter potential: $V_{P,e}=0.2\,V_{CC}$ changes $P(\nu_\mu\to\nu_e)$ at DUNE from 0.0723269 to 0.0770500.

\Needspace{12\baselineskip}\section*{12. Verification}\addcontentsline{toc}{section}{12. Verification}
Every statement above is a check of 7+1.nb (Wolfram Language 15.0.1) and of the independent Python program tools/ref\_sm.py (numpy, scipy, sympy). The numbers printed by both programs:

{\footnotesize
\begin{longtable}{>{\raggedright\arraybackslash}p{\dimexpr 0.4200\linewidth-2\tabcolsep\relax}>{\raggedright\arraybackslash}p{\dimexpr 0.1000\linewidth-2\tabcolsep\relax}>{\raggedright\arraybackslash}p{\dimexpr 0.2000\linewidth-2\tabcolsep\relax}>{\raggedright\arraybackslash}p{\dimexpr 0.2800\linewidth-2\tabcolsep\relax}}
\toprule
\textbf{quantity} & \textbf{values} & \textbf{max. relative difference} & \textbf{note} \\
\midrule\endfirsthead
\toprule
\textbf{quantity} & \textbf{values} & \textbf{max. relative difference} & \textbf{note} \\
\midrule\endhead
\endfoot
\bottomrule\endlastfoot
$v,g,g',\sin^2\theta_W,e,1/\alpha,\lambda$ & 7 & $2.1\times10^{-16}$ &  \\
$|V_{\rm CKM}|$ & 9 & $2.3\times10^{-5}$ & Python prints 6 decimals \\
$J_{\rm CKM}$ & 1 & $2.2\times10^{-16}$ &  \\
$\bar\rho,\bar\eta$ & 2 & $3.5\times10^{-16}$ &  \\
$|U_{\rm PMNS}|$ & 9 & $3.2\times10^{-5}$ & Python prints 5 decimals \\
$J_{\rm PMNS}$ & 1 & $1.2\times10^{-16}$ &  \\
$m_1,m_2,m_3$ (eV) & 3 & 0 &  \\
$y_t$ & 1 & $1.1\times10^{-16}$ &  \\
age (Gyr) & 1 & $1.0\times10^{-15}$ &  \\
$V_\perp$, $V_\perp^{1/4}$, lowest KK mass & 3 & $2.2\times10^{-16}$ &  \\
light (eV) and heavy (GeV) masses with the portal & 7 & $9.9\times10^{-7}$ &  \\
$(1-NN^\dagger)_{ee}$ & 1 & $3.3\times10^{-7}$ &  \\
DUNE vacuum $P(\mu\to e)$, $P(\bar\mu\to\bar e)$, $P(\mu\to\mu)$ & 3 & $4.3\times10^{-16}$ &  \\
DUNE matter $P(\mu\to e)$, $P(\bar\mu\to\bar e)$ & 2 & $1.2\times10^{-15}$ &  \\
vector portal: $P(\mu\to e)$ with, without & 2 & $9.0\times10^{-16}$ &  \\
Sun $P_{ee}$ ($E$, numerical, adiabatic) & 12 & $3.8\times10^{-6}$ &  \\
supernova resonance densities H, L & 2 & $1.9\times10^{-16}$ &  \\
$T_\nu$, $n_\nu$ & 2 & $1.3\times10^{-16}$ &  \\*
\end{longtable}}
Apart from the matrices that the Python program prints rounded, the largest relative difference is $3.8\times10^{-6}$.

The checks of sections 17–19 of the notebook:

{\footnotesize
\begin{longtable}{>{\raggedright\arraybackslash}p{\dimexpr 0.0600\linewidth-2\tabcolsep\relax}>{\raggedright\arraybackslash}p{\dimexpr 0.7400\linewidth-2\tabcolsep\relax}>{\raggedright\arraybackslash}p{\dimexpr 0.0800\linewidth-2\tabcolsep\relax}>{\raggedright\arraybackslash}p{\dimexpr 0.1200\linewidth-2\tabcolsep\relax}}
\toprule
\textbf{\#} & \textbf{check} & \textbf{result} & \textbf{section} \\
\midrule\endfirsthead
\toprule
\textbf{\#} & \textbf{check} & \textbf{result} & \textbf{section} \\
\midrule\endhead
\endfoot
\bottomrule\endlastfoot
1 & \texttt{brane frame metric diag(+1, -1, -1, -1) == eta17 restricted to x0..x3: the squares of the octonion basis e[0], e[1], e[2], e[3] of the quaternion subalgebra (Re(e[A] e[B]), A, B = 0..3), which is closed under the octonion product} & True & 17. \\
2 & \texttt{anomaly cancellation per generation (with nu\_R): SU(3)\textasciicircum{}2U(1), SU(2)\textasciicircum{}2U(1), U(1)\textasciicircum{}3, grav\textasciicircum{}2U(1), SU(3)\textasciicircum{}3 all vanish; doublets per generation even} & True & 17.1 \\
3 & \texttt{electric charges Q = T3 + Y: (u, d) = (2/3, -1/3), (nu, e) = (0, -1), u\_R 2/3, d\_R -1/3, e\_R -1, nu\_R 0} & True & 17.1 \\
4 & \texttt{g* = 28 + (7/8) 90 = 106.75 for the Standard Model fields (without nu\_R)} & True & 17.1 \\
5 & \texttt{gauge-boson mass matrix: eigenvalues g\textasciicircum{}2 v\textasciicircum{}2/4 (twice), (g\textasciicircum{}2 + g'\textasciicircum{}2) v\textasciicircum{}2/4 and 0 (photon)} & True & 17.3 \\
6 & \texttt{tree level: m\_W == g v/2, m\_Z == Sqrt[g\textasciicircum{}2 + g'\textasciicircum{}2] v/2, rho == m\_W\textasciicircum{}2/(m\_Z\textasciicircum{}2 cos\textasciicircum{}2 theta\_W) == 1, m\_H\textasciicircum{}2 == 2 lambda v\textasciicircum{}2} & True & 17.3 \\
7 & \texttt{CKM: unitary; Jarlskog invariant J = 3.12 x 10\textasciicircum{}-5 (PDG 2024: 3.12 +0.13 -0.12 x 10\textasciicircum{}-5); A, rhobar, etabar within the PDG uncertainties} & True & 17.4 \\
8 & \texttt{Yukawa matrices: singular values x v/Sqrt[2] are the quark masses and |U\_uL\textasciicircum{}dag.U\_dL| == |V\_CKM|} & True & 17.4 \\
9 & \texttt{PMNS unitary; Dirac neutrino Yukawa Sqrt[2] U.diag(m)/v reproduces the masses (singular values)} & True & 17.4 \\
10 & \texttt{age of the universe from the brane Friedmann equation: 13.79 Gyr (Planck 2018: 13.787 +- 0.020 Gyr)} & True & 17.5 \\
11 & \texttt{ADD relation: kappa8 rho4/V\_perp == 8 Pi G4 rho4 with G4 = G8/V\_perp (kappa8 = 8 Pi G8)} & True & 17.5 \\
12 & \texttt{brane representation: R\_A R\_B + R\_B R\_A == 2 eta17\_AB ID16 for R = (gamma\_mu (x) 1\_4, gamma5 (x) rho\_k)} & True & 18. \\
13 & \texttt{intertwiner: Gamma17[A].U == U.R\_A for A = 0..7, and U is unitary} & True & 18. \\
14 & \texttt{chi17 == U.(gamma5 (x) rho4567).U\textasciicircum{}dag; rho4567 Hermitian with eigenvalues \{-1, -1, 1, 1\}} & True & 18. \\
15 & \texttt{bilinears at the brane: psibar psi' == Sum\_a psibar\_a psi'\_a and psibar Gamma\textasciicircum{}mu psi' == Sum\_a psibar\_a gamma\textasciicircum{}mu psi'\_a (mu = 0..3)} & True & 18. \\
16 & \texttt{KK operator: X.X == -|q|\textasciicircum{}2/b0\textasciicircum{}2 ID16 and X anticommutes with gamma\_mu (x) 1\_4} & True & 18.1 \\
17 & \texttt{KK masses: Det[gamma.p (x) 1 + X - m] vanishes on p0\textasciicircum{}2 - |p|\textasciicircum{}2 == m\textasciicircum{}2 + |q|\textasciicircum{}2/b0\textasciicircum{}2 (40-digit arithmetic) and not off shell} & True & 18.1 \\
18 & \texttt{portal mixing: 1 - N.N\textasciicircum{}dag == mu.mu\textasciicircum{}dag/m\textasciicircum{}2 to second order (non-unitary light mixing)} & True & 18.2 \\
19 & \texttt{vacuum: P(a->b) doubly stochastic; equals delta - 4 Sum Re X sin\textasciicircum{}2 Delta + 2 Sum Im X sin 2Delta; CPT P(a->b) == P(bbar->abar)} & True & 19. \\
20 & \texttt{V\_CC(rho Y\_e = 1 g/cm\textasciicircum{}3) == 7.63 x 10\textasciicircum{}-14 eV; matter evolution doubly stochastic; a flavour-universal potential has no effect} & True & 19. \\
21 & \texttt{two-flavour MSW: sin\textasciicircum{}2 theta\_m from diagonalization == formula, and sin\textasciicircum{}2 2theta\_m == 1 at 2 E V == dm\textasciicircum{}2 cos 2theta} & True & 19. \\
22 & \texttt{Sun: numerical P\_ee == adiabatic MSW formula to 2 x 10\textasciicircum{}-3 at 0.3, 1, 5 and 10 MeV; P\_ee(0.3 MeV) within 0.03 of the vacuum average; P\_ee falls with energy} & True & 19. \\
23 & \texttt{cosmic neutrino background: T\_nu == 1.945 K and n == 112 cm\textasciicircum{}-3 per flavour} & True & 19. \\
24 & \texttt{dirac16complex portal: non-unitary oscillations, Sum\_b P(a->b) < 1 (zero-distance effect); decoupled limit == three-flavour vacuum result} & True & 19. \\*
\end{longtable}}
The Python program’s boolean checks:

{\footnotesize
\begin{longtable}{>{\raggedright\arraybackslash}p{\dimexpr 0.0600\linewidth-2\tabcolsep\relax}>{\raggedright\arraybackslash}p{\dimexpr 0.8200\linewidth-2\tabcolsep\relax}>{\raggedright\arraybackslash}p{\dimexpr 0.1200\linewidth-2\tabcolsep\relax}}
\toprule
\textbf{\#} & \textbf{check} & \textbf{result} \\
\midrule\endfirsthead
\toprule
\textbf{\#} & \textbf{check} & \textbf{result} \\
\midrule\endhead
\endfoot
\bottomrule\endlastfoot
1 & \texttt{anomaly cancellation in each generation (with nu\_R): SU(3)\textasciicircum{}2U(1), SU(2)\textasciicircum{}2U(1), U(1)\textasciicircum{}3, grav\textasciicircum{}2U(1), SU(3)\textasciicircum{}3 all 0} & True \\
2 & \texttt{Witten SU(2) anomaly: number of doublets per generation (3 colours of Q + 1 L) = 4, even} & True \\
3 & \texttt{electric charges Q = T3 + Y: Q -> (2/3, -1/3), u\_R 2/3, d\_R -1/3, L -> (0, -1), e\_R -1, nu\_R 0} & True \\
4 & \texttt{effective number of relativistic degrees of freedom g* = 28 + (7/8) 90 = 106.75 (106.75 + 5.25 if nu\_R were thermal)} & True \\
5 & \texttt{gauge-boson mass matrix from |D<Phi>|\textasciicircum{}2: eigenvalues g\textasciicircum{}2v\textasciicircum{}2/4 (twice, W+-), (g\textasciicircum{}2+g'\textasciicircum{}2)v\textasciicircum{}2/4 (Z), 0 (photon)} & True \\
6 & \texttt{tree-level relations: m\_W = g v/2, m\_Z = sqrt(g\textasciicircum{}2+g'\textasciicircum{}2) v/2, rho = m\_W\textasciicircum{}2/(m\_Z\textasciicircum{}2 cos\textasciicircum{}2) = 1, m\_H\textasciicircum{}2 = 2 lambda v\textasciicircum{}2} & True \\
7 & \texttt{CKM unitarity V V\textasciicircum{}dag = 1} & True \\
8 & \texttt{J = c12 c23 c13\textasciicircum{}2 s12 s23 s13 sin(delta) = 3.12e-5 (PDG 2024: 3.12 +0.13 -0.12 e-5)} & True \\
9 & \texttt{Wolfenstein parameters agree with the PDG fit (A 0.826, rhobar 0.1591, etabar 0.3523) within their quoted uncertainties} & True \\
10 & \texttt{Yukawa matrices Y\_u = sqrt2 diag(m\_u,m\_c,m\_t)/v, Y\_d = sqrt2 V diag(m\_d,m\_s,m\_b)/v: singular values x v/sqrt2 = masses, |U\_uL\textasciicircum{}dag U\_dL| = |V|} & True \\
11 & \texttt{PMNS unitarity} & True \\
12 & \texttt{neutrino Yukawa matrix Y\_nu = sqrt2 U diag(m1,m2,m3)/v: masses and |U| recovered (singular value decomposition)} & True \\
13 & \texttt{vacuum: probabilities P(a -> b) form a doubly stochastic matrix (rows and columns sum to 1)} & True \\
14 & \texttt{vacuum: amplitude result = delta\_ab - 4 sum\_\{i>j\} Re(X\_ij) sin\textasciicircum{}2 Delta\_ij + 2 sum\_\{i>j\} Im(X\_ij) sin 2Delta\_ij, X\_ij = U\_ai* U\_bi U\_aj U\_bj*} & True \\
15 & \texttt{CPT: P(a->b) = P(bbar->abar); CP asymmetry P(mu->e) - P(mubar->ebar) proportional to J (vanishes for delta = 0, pi)} & True \\
16 & \texttt{V\_CC(rho Y\_e = 1 g/cm\textasciicircum{}3) = 7.63e-14 eV} & True \\
17 & \texttt{matter (constant density): unitary evolution, probabilities doubly stochastic} & True \\
18 & \texttt{matter: a flavour-universal potential has no effect (only differences of potentials matter)} & True \\
19 & \texttt{two-flavour MSW: sin\textasciicircum{}2 theta\_m from diagonalization = formula; maximal mixing (sin\textasciicircum{}2 2theta\_m = 1) at 2 E V = dm\textasciicircum{}2 cos 2theta} & True \\
20 & \texttt{Sun: numerical P\_ee = adiabatic MSW formula to 2e-3 at E = 0.3, 1, 5, 10 MeV} & True \\
21 & \texttt{Sun: P\_ee falls from near the vacuum average c13\textasciicircum{}4 (1 - sin\textasciicircum{}2(2 theta12)/2) + s13\textasciicircum{}4 (within 0.03 at 0.3 MeV) to about 0.31 at 10 MeV} & True \\
22 & \texttt{CNB: T\_nu = 1.945 K and n = 112 cm\textasciicircum{}-3 per flavour} & True \\
23 & \texttt{brane frame metric diag(+1,-1,-1,-1) = eta17 on x0..x3 = squares of the quaternion basis 1, e\_1, e\_2, e\_3 (Re(e\_a e\_b)); H closed} & True \\
24 & \texttt{brane representation: gamma\_mu (x) 1\_4 (mu = 0..3) and gamma5 (x) rho\_k (k = 4..7) satisfy \{R\_A, R\_B\} = 2 eta\_AB} & True \\
25 & \texttt{intertwiner U (Schur average over the 256 Clifford products): Gamma17[A] U = U R\_A for A = 0..7, U unitary} & True \\
26 & \texttt{8D chirality at the brane: chi17 = U (gamma5 (x) rho4567) U\textasciicircum{}dag, rho4567 Hermitian with eigenvalues +1, +1, -1, -1} & True \\
27 & \texttt{bilinears at the brane: Psibar Psi' = sum\_a psibar\_a psi'\_a and Psibar Gamma\textasciicircum{}mu Psi' = sum\_a psibar\_a gamma\textasciicircum{}mu psi'\_a (mu = 0..3)} & True \\
28 & \texttt{Kaluza-Klein: the 16x16 operator gamma.p (x) 1 + gamma5 (x) rho.q/b0 - m is singular exactly on p0\textasciicircum{}2 - |p|\textasciicircum{}2 = m\textasciicircum{}2 + |q|\textasciicircum{}2/b0\textasciicircum{}2} & True \\
29 & \texttt{KK operator: X\textasciicircum{}2 = -|q|\textasciicircum{}2/b0\textasciicircum{}2 and X anticommutes with gamma\textasciicircum{}mu (x) 1} & True \\
30 & \texttt{non-unitary oscillations: sum\_b P(a->b) = (N N\textasciicircum{}dag)\_aa\textasciicircum{}2-type deficit < 1 (zero-distance effect), complemented by the heavy states} & True \\
31 & \texttt{portal decoupling: for mu -> 0 the light block becomes the PMNS matrix and the oscillation probabilities the three-flavour vacuum ones (section 19)} & True \\
32 & \texttt{age 13.79 Gyr (Planck 2018: 13.787 +- 0.020)} & True \\
33 & \texttt{ADD relation: 3 Ha\textasciicircum{}2 = kappa8 rho4/V\_perp = 8 pi G4 rho4 with G4 = G8/V\_perp} & True \\*
\end{longtable}}
\Needspace{12\baselineskip}\section*{13. How to reproduce}\addcontentsline{toc}{section}{13. How to reproduce}
\begin{enumerate}
\item \texttt{python3 tools/make\_7p1\_nb.py} writes 7+1.nb (unevaluated).
\item Evaluate 7+1.nb in a fresh kernel of Wolfram 15.0.1 (Evaluation \texorpdfstring{\ensuremath{\triangleright}}{▸} Quit Kernel, then Evaluation \texorpdfstring{\ensuremath{\triangleright}}{▸} Evaluate Notebook; or \texttt{Get["D:/Developer/github/7+1/tools/run\_eval.wl"]}) and save it; the last section summarizes all checks.
\item \texttt{python3 tools/ref\_sm.py} repeats the checks of this document independently (exit code 0 if all pass).
\item \texttt{python3 tools/make\_docs.py} regenerates this document, its symbol dictionary and its provenance record.
\end{enumerate}
\Needspace{12\baselineskip}\section*{14. References}\addcontentsline{toc}{section}{14. References}
\begin{itemize}
\item S. Navas et al. (Particle Data Group), \emph{Review of Particle Physics}, Phys. Rev. D 110 (2024) 030001, and its 2025 update (pdg.lbl.gov) — $G_F$, $m_W$, $m_Z$, $m_H$, quark masses, the CKM parameters and $J$.
\item I. Esteban, M. C. Gonzalez-Garcia, M. Maltoni, I. Martinez-Soler, J. P. Pinheiro, T. Schwetz, \emph{NuFit-6.0: updated global analysis of three-flavor neutrino oscillations}, JHEP 12 (2024) 216, arXiv:2410.05380 — the PMNS parameters.
\item Planck Collaboration, N. Aghanim et al., \emph{Planck 2018 results. VI. Cosmological parameters}, Astron. Astrophys. 641 (2020) A6.
\item N. Arkani-Hamed, S. Dimopoulos, G. Dvali, Phys. Lett. B 429 (1998) 263 — gravity in large extra dimensions, $G_4=G_8/V_\perp$.
\item L. Wolfenstein, Phys. Rev. D 17 (1978) 2369; S. P. Mikheyev, A. Yu. Smirnov, Sov. J. Nucl. Phys. 42 (1985) 913 — the MSW effect.
\item J. N. Bahcall, \emph{Neutrino Astrophysics}, Cambridge University Press (1989) — the exponential solar electron density.
\item C. Jarlskog, Phys. Rev. Lett. 55 (1985) 1039 — the invariant $J$.
\item E. Witten, Phys. Lett. B 117 (1982) 324 — the global $SU(2)$ anomaly.
\end{itemize}
\end{document}
